% GENERATED FILE. Edit canonical lessons, then run npm run generate:books.
% canonical-source-hash: fnv1a64:d1e6a1945711f38a
% canonical-lessons: GE-C277-streik, GE-C277-streiken, GE-C277-protestieren, GE-C277-demonstrieren, GE-C277-krise, GE-R277-review-270-277

\chapter{Der Streik}
\label{ch:ge-der-streik}
\hlchaptermodality{277}

\begin{chapteropening}
\textbf{By the end of this chapter:} \emph{I can report protests: a strike, to go on strike, to protest, to demonstrate and a crisis.}

Complete the last lesson of chapter 277: I can report protests: a strike, to go on strike, to protest, to demonstrate and a crisis.
\end{chapteropening}

\section[der Streik]{der Streik — the strike (at work)}

\label{lesson:GE-C277-streik}



\begin{quote}
Before the new one: say the German for to increase, then the German for to sink.
\end{quote}

\subsection*{You'll want to know: der Streik}
\textbf{der Streik} — \textquotedblleft{}the strike (at work)\textquotedblright{}.

\begin{quote}
Wegen des Streiks fuhren keine Züge.
\end{quote}

\begin{quote}
\emph{Because of the strike no trains ran.}
\end{quote}

Plural: \emph{die Streiks}.

\begin{cousinweb}
\emph{Streik} was borrowed from English \emph{strike} (EtymDB) and given German spelling.
\end{cousinweb}

\begin{grammarlens}[title={What you've built}]
One more word for telling, explaining and describing at B1.
\end{grammarlens}

\subsection*{Your turn}
\begin{itemize}
\raggedright
  \item \emph{Say it:} \emph{der Streik}
  \item \emph{Say it:} \emph{der Streik}, once more
  \item \emph{Say it:} a sentence of your own with \emph{der Streik}
  \item \emph{Recall:} say the German for to increase, then the German for to sink, then say \emph{der Streik} again
\end{itemize}

\begin{tcolorbox}[breakable,colback=teal!4,colframe=teal!35!black,title={Before you move on}]
What does der Streik mean? (\textquotedblleft{}The strike (at work)\textquotedblright{}.) Say it once more.
\end{tcolorbox}
\section[streiken]{streiken — to go on strike}

\label{lesson:GE-C277-streiken}



\begin{quote}
Before the new one: say the German for to sink, then the German for the strike.
\end{quote}

\subsection*{You'll want to know: streiken}
\textbf{streiken} — \textquotedblleft{}to go on strike\textquotedblright{}.

\begin{quote}
Die Busfahrer streikten zwei Tage lang.
\end{quote}

\begin{quote}
\emph{The bus drivers were on strike for two days.}
\end{quote}

Regular: \emph{sie streikt}, \emph{sie hat gestreikt}.

\begin{cousinweb}
\emph{streiken} is built on \emph{Streik}, which was borrowed from English \emph{strike} (EtymDB).
\end{cousinweb}

\begin{grammarlens}[title={What you've built}]
One more word for telling, explaining and describing at B1.
\end{grammarlens}

\subsection*{Your turn}
\begin{itemize}
\raggedright
  \item \emph{Say it:} \emph{streiken}
  \item \emph{Say it:} \emph{streiken}, once more
  \item \emph{Say it:} a sentence of your own with \emph{streiken}
  \item \emph{Recall:} say the German for to sink, then the German for the strike, then say \emph{streiken} again
\end{itemize}

\begin{tcolorbox}[breakable,colback=teal!4,colframe=teal!35!black,title={Before you move on}]
What does streiken mean? (\textquotedblleft{}To go on strike\textquotedblright{}.) Say it once more.
\end{tcolorbox}
\section[protestieren]{protestieren — to protest}

\label{lesson:GE-C277-protestieren}



\begin{quote}
Before the new one: say the German for the strike, then the German for to go on strike.
\end{quote}

\subsection*{You'll want to know: protestieren}
\textbf{protestieren} — \textquotedblleft{}to protest\textquotedblright{}.

\begin{quote}
Viele Bürger protestierten gegen die neue Straße.
\end{quote}

\begin{quote}
\emph{Many citizens protested against the new road.}
\end{quote}

With \emph{gegen}: \emph{gegen etwas protestieren}. Perfekt without \emph{ge-}: \emph{hat protestiert}.

\begin{cousinweb}
\emph{protestieren} was borrowed from French \emph{protester}, from Latin \emph{prōtestārī} (EtymDB) \textquotedblleft{}to declare in public\textquotedblright{}, built on \emph{testis} \textquotedblleft{}a witness\textquotedblright{} (Lewis \& Short).
\end{cousinweb}

\begin{grammarlens}[title={What you've built}]
One more word for telling, explaining and describing at B1.
\end{grammarlens}

\subsection*{Your turn}
\begin{itemize}
\raggedright
  \item \emph{Say it:} \emph{protestieren}
  \item \emph{Say it:} \emph{protestieren}, once more
  \item \emph{Say it:} a sentence of your own with \emph{protestieren}
  \item \emph{Recall:} say the German for the strike, then the German for to go on strike, then say \emph{protestieren} again
\end{itemize}

\begin{tcolorbox}[breakable,colback=teal!4,colframe=teal!35!black,title={Before you move on}]
What does protestieren mean? (\textquotedblleft{}To protest\textquotedblright{}.) Say it once more.
\end{tcolorbox}
\section[demonstrieren]{demonstrieren — to demonstrate}

\label{lesson:GE-C277-demonstrieren}



\begin{quote}
Before the new one: say the German for to go on strike, then the German for to protest.
\end{quote}

\subsection*{You'll want to know: demonstrieren}
\textbf{demonstrieren} — \textquotedblleft{}to demonstrate\textquotedblright{}.

\begin{quote}
Am Samstag demonstrierten tausend Menschen in der Stadt.
\end{quote}

\begin{quote}
\emph{On Saturday a thousand people demonstrated in the town.}
\end{quote}

Short noun: \emph{die Demo}. Also to show how something works.

\begin{cousinweb}
\emph{demonstrieren} is a learned borrowing from Latin \emph{dēmonstrāre} \textquotedblleft{}to show\textquotedblright{} (EtymDB), built on \emph{monstrāre} \textquotedblleft{}to show\textquotedblright{} (Lewis \& Short).
\end{cousinweb}

\begin{grammarlens}[title={What you've built}]
One more word for telling, explaining and describing at B1.
\end{grammarlens}

\subsection*{Your turn}
\begin{itemize}
\raggedright
  \item \emph{Say it:} \emph{demonstrieren}
  \item \emph{Say it:} \emph{demonstrieren}, once more
  \item \emph{Say it:} a sentence of your own with \emph{demonstrieren}
  \item \emph{Recall:} say the German for to go on strike, then the German for to protest, then say \emph{demonstrieren} again
\end{itemize}

\begin{tcolorbox}[breakable,colback=teal!4,colframe=teal!35!black,title={Before you move on}]
What does demonstrieren mean? (\textquotedblleft{}To demonstrate\textquotedblright{}.) Say it once more.
\end{tcolorbox}
\section[die Krise]{die Krise — the crisis}

\label{lesson:GE-C277-krise}



\begin{quote}
Before the new one: say the German for to protest, then the German for to demonstrate.
\end{quote}

\subsection*{You'll want to know: die Krise}
\textbf{die Krise} — \textquotedblleft{}the crisis\textquotedblright{}.

\begin{quote}
Das Land war in einer tiefen Krise.
\end{quote}

\begin{quote}
\emph{The country was in a deep crisis.}
\end{quote}

Plural: \emph{die Krisen}.

\begin{cousinweb}
\emph{Krise} was borrowed through Latin \emph{crisis} from Greek \emph{krisis} \textquotedblleft{}a decision, the turning point of an illness\textquotedblright{}; EtymDB lists no parent, so this follows the usual account.
\end{cousinweb}

\begin{grammarlens}[title={What you've built}]
One more word for telling, explaining and describing at B1.
\end{grammarlens}

\subsection*{Your turn}
\begin{itemize}
\raggedright
  \item \emph{Say it:} \emph{die Krise}
  \item \emph{Say it:} \emph{die Krise}, once more
  \item \emph{Say it:} a sentence of your own with \emph{die Krise}
  \item \emph{Recall:} say the German for to protest, then the German for to demonstrate, then say \emph{die Krise} again
\end{itemize}

\begin{tcolorbox}[breakable,colback=teal!4,colframe=teal!35!black,title={Before you move on}]
What does die Krise mean? (\textquotedblleft{}The crisis\textquotedblright{}.) Say it once more.
\end{tcolorbox}
\section[(dialogue)]{Words from chapters 270–277 — a review}

\label{lesson:GE-R277-review-270-277}



\begin{quote}
Say the German for to hesitate and to wait and see.
\end{quote}

\subsection*{Your turn}
Say these aloud:

\begin{itemize}
\raggedright
  \item to hesitate, to wait and see, to delay, to go, the sequence of events
  \item to report, the news item, the report, the headline, the article
  \item the broadcast, radio, the press, the journalist, the front page
  \item to publish, to mention, to express, to announce, to make public
  \item the rumour, the source, to claim, the claim, to confirm
  \item supposedly, apparently, according to, to suspect, presumably
  \item the survey, the statistics, the share, to increase, to sink
  \item the strike, to go on strike, to protest, to demonstrate, the crisis
\end{itemize}

\begin{tcolorbox}[breakable,colback=teal!4,colframe=teal!35!black,title={Before you move on}]
To hesitate? (\textbf{zögern}.) To wait and see? (\textbf{abwarten}.) To delay? (\textbf{verzögern}.) To go? (\textbf{verlaufen}.) The sequence of events? (\textbf{der Ablauf}.) To report? (\textbf{melden}.) The news item? (\textbf{die Meldung}.) The report? (\textbf{der Bericht}.) The headline? (\textbf{die Schlagzeile}.) The article? (\textbf{der Artikel}.) The broadcast? (\textbf{die Sendung}.) Radio? (\textbf{der Rundfunk}.) The press? (\textbf{die Presse}.) The journalist? (\textbf{der Journalist}.) The front page? (\textbf{die Titelseite}.) To publish? (\textbf{veröffentlichen}.) To mention? (\textbf{erwähnen}.) To express? (\textbf{äußern}.) To announce? (\textbf{ankündigen}.) To make public? (\textbf{bekanntgeben}.) The rumour? (\textbf{das Gerücht}.) The source? (\textbf{die Quelle}.) To claim? (\textbf{behaupten}.) The claim? (\textbf{die Behauptung}.) To confirm? (\textbf{bestätigen}.) Supposedly? (\textbf{angeblich}.) Apparently? (\textbf{offenbar}.) According to? (\textbf{zufolge}.) To suspect? (\textbf{vermuten}.) Presumably? (\textbf{vermutlich}.) The survey? (\textbf{die Umfrage}.) The statistics? (\textbf{die Statistik}.) The share? (\textbf{der Anteil}.) To increase? (\textbf{zunehmen}.) To sink? (\textbf{sinken}.) The strike? (\textbf{der Streik}.) To go on strike? (\textbf{streiken}.) To protest? (\textbf{protestieren}.) To demonstrate? (\textbf{demonstrieren}.) The crisis? (\textbf{die Krise}.)
\end{tcolorbox}
